\appendix
\section{Computational and numerical reproducibility}\label{sec:computing}
The GPU experiment used Python 3.13.13, PyTorch 2.14.0 with CUDA 13.0, NeuralForecast 3.2.2, the TimesFM package 3.0.2, Chronos-forecasting 2.3.2 and scikit-learn 1.9.0. The pretrained checkpoint revisions were \texttt{1d952420fba87f3c6dee4f240de0f1a0fbc790e3} for TimesFM 2.5 and \texttt{29ec3766d36d6f73f0696f85560a422f50e8498c} for Chronos-2. TimesFM received 16 observed context records, padded to a maximum context of 32, with a configured maximum horizon of 32; only its first forecast position was evaluated. Chronos used float32, disabled cross-query learning and returned the median forecast. Model configurations, checkpoints and record-level outputs are retained in the repository; pretrained weights are identified by revision rather than redistributed.

All 128,896 logged donor relationships passed the checks for same-pelleting-line membership, partition membership and completion before the target origin; the total includes repeated relationships across nested and expanding blocks. Source-data, frozen-model and final-prediction fingerprints remained unchanged, and the shared CPU and GPU data, window and encoding functions were checked for consistency. These checks concern the recorded completion-time proxy, since actual publication delays are not available in the source.

\fig{12_numerical_and_seeds}{numerics}{1}{Numerical and training variability. (a) Every paired CPU/GPU Chronos-2 covariate output as a signed difference against the CPU prediction. (b) All three GPU seeds (11, 29 and 47) and their prediction-mean ensemble for TFT and N-HiTS, with the matched Ridge reference; averaging predictions is not the same as averaging the three MAEs. These diagnostics accompany the primary comparison without selecting a preferred seed or device.}

Within the GPU run, the 24 saved neural checkpoints were reloaded and compared with their stored inference, together with four Chronos covariate batch-size checks. All 28 checks passed the recorded tolerances with a maximum difference of $7.63\times10^{-5}$ kWh. Fresh CPU and GPU neural fits were less closely aligned, as optimization can follow different trajectories even with the same integer seed: TFT's CPU and GPU MAEs were 53.25 and 53.65 kWh, and neither established an advantage over the matched Ridge control under the paired intervals. Figure~\ref{fig:numerics} shows all three GPU seeds and their ensemble.

For frozen Chronos-2 with covariates, the median absolute CPU--GPU prediction difference was 0.86 kWh and the maximum 57.04 kWh, and its overall MAE moved from 85.65 to 85.53 kWh. Maximum prediction differences for both frozen TimesFM variants and for univariate Chronos stayed below $1.9\times10^{-4}$ kWh. To trace the covariate sensitivity, the installed Chronos preprocessing and instance normalization were applied to the same 584 contexts, giving 6,424 target/covariate channels. Among the 1,097 constant historical channels, 694 showed a CPU--GPU difference above 0.1 in normalized context values, whereas no non-constant channel exceeded that threshold; the largest difference was 0.88137 in the transformed coordinates.

The implementation substitutes a small scale only when the calculated scale is exactly zero, so tiny round-off variances in constant channels can be amplified before the inverse-hyperbolic-sine transform. This mechanism is consistent with the observed sensitivity, although it is not an end-to-end causal demonstration for each prediction. The reported benchmark keeps the native preprocessing and original outputs; a corrective normalization policy would require a separately selected and evaluated experiment.
